\documentclass[12pt]{article}

\usepackage{bookmark}
\usepackage{subfiles}
\usepackage[explicit, newparttoc]{titlesec}
\usepackage[]{titletoc}
\usepackage[titles]{tocloft}
\usepackage{enumitem}
\usepackage[most]{tcolorbox}
\usepackage[]{lipsum}
\usepackage{palatino}
\usepackage{xltabular}
\usepackage{csquotes}
\usepackage{tikz}

\usepackage[a4paper,left=2cm,right=2cm,top=3cm,bottom=2cm]{geometry}

% Set paragraphs to have vertical skip and no indent
\usepackage[parfill]{parskip}

\usetikzlibrary{calc}

\bookmarksetup{numbered}
\graphicspath{{resources/}}

%% --------------------------- Section, Division, Part Numbering and Format
\counterwithout{subsection}{section}
\counterwithin*{section}{part}
\renewcommand*{\thepart}{\arabic{part}}

\titleformat{\section}[hang]{\normalfont\Large\bfseries}{Division \thesection\ - }{0em}{#1}
\titleformat{\part}[hang]{\normalfont\Large\bfseries}{Part \thepart\ - }{0em}{#1}
\titleformat{name = \part, numberless}[hang]{\normalfont\Large\bfseries}{}{0em}{#1}

\titlecontents{part}
  [0pt]
  {\addvspace{1pc}}%
  {\normalfont\large\bfseries PART \thecontentslabel\ - }
  {\bfseries \large PART \contentsmargin{0pt}}
  {}
  [\addvspace{.5pc}]


% new section for non numbered parts
\newcounter{nonnumberedparts}
\titleclass{\nonnumberedparts}{straight}[\part]
\titleformat{\nonnumberedparts}[hang]{\normalfont\Large\bfseries}{}{0em}{#1}
\titlespacing*{\nonnumberedparts}{0pt}{3.5ex plus 1ex minus .2ex}{2.3 ex plus .2ex}
\renewcommand{\thenonnumberedparts}{\arabic{nonnumberedparts}}
% \counterwithin{section}{mysectionnonnumberedparts}

% define appearance for new section in toc
\titlecontents{nonnumberedparts}
  [0pt]
  {\addvspace{1pc}}%
  {\normalfont\large\bfseries}
  {\bfseries \large \contentsmargin{0pt}}
  {}
  [\addvspace{.5pc}]

%make formatting 

%% Correct 'Division' labelling in TOC
\renewcommand{\cftsecpresnum}{Division }
\renewcommand{\cftsecaftersnum}{ --- }
\newlength{\mylen} % a "scratch" length
\settowidth{\mylen}{\bfseries\cftsecpresnum\cftsecaftersnum} % extra space
\addtolength{\cftsecnumwidth}{\mylen} % add the extra space

% \makeatletter
% \def\@part[#1]#2{%
%     \ifnum \c@secnumdepth >-2\relax
%       \refstepcounter{part}%
%       \addcontentsline{toc}{part}{\partname~\thepart\hspace{1em}#1}%
%     \else
%       \addcontentsline{toc}{part}{#1}%
%     \fi
%     \markboth{}{}%
%     {\centering
%      \interlinepenalty \@M
%      \normalfont
%      \ifnum \c@secnumdepth >-2\relax
%        \huge\bfseries \partname\nobreakspace\thepart
%        \par
%        \vskip 20\p@
%      \fi
%      \Huge \bfseries #2\par}%
%     \@endpart}
% \makeatother

% --------------------- Numbered Items
\setlist{leftmargin=*}
\setlist[enumerate,1]{label = \thesubsection.\arabic*., 
                        ref = \thesubsection.\arabic*,
                        leftmargin = 28pt}
\setlist[enumerate,2]{label = \theenumi.\alph*.,
                        ref = \theenumi.\alph*,
                        leftmargin = 28pt}
\setlist[enumerate,3]{label = \theenumii.\roman*.,     
                        ref = \theenumii.\roman*,
                        leftmargin = 32pt}
\setlist[enumerate,4]{label = \theenumiii.(\arabic*).,     
                        ref = \theenumiii.(\arabic*),
                        leftmargin = 40pt}

\setlist[itemize,1]{label=-}

% --------------------- Flags

\newif\ifshownotes

% to toggle on and off the notes blocks, comment and uncomment
% the shownotestrue and shownotesfalse below
\shownotestrue
% \shownotesfalse

% \ifshownotes
%  % only shown if \importanttrue is set
%  this is some important text
% \fi

% --------------------- Notes Block

\newtcolorbox{noteblock}[1][]{
  title=#1,
  % nobeforeafter,
  enhanced,
  % sharp corners,
  % boxrule=0pt,
  % frame hidden,
  toprule=1pt,
  bottomrule=1pt,
  leftrule=1pt,
  rightrule=1pt,
  colback=blue!5!white,
  colbacktitle=blue,
  colframe=blue
}


\makeatletter
\NewDocumentCommand\note{+O{Note}+m}{%
\ifshownotes
 \begingroup
% \tcbset{%
%   noteshift/.store in=\mynote@shift,
%  noteshift=1.5cm
% }
  \begin{noteblock}[#1]
   #2
 \end{noteblock}
 \endgroup
 \fi
 }
\makeatother

\makeatletter
% Parskip in minipages (e.g. tcolorboxes.)
\newcommand{\@minipagerestore}{\setlength{\parskip}{\medskipamount}}
\makeatother

% Headers and footers
\usepackage{fancyhdr}
\fancyhead{}
\fancyhead[R]{Constitution of Perth Artifactory Inc}
\fancyhead[L]{%
\begin{tikzpicture}[overlay, remember picture]%
    \node[anchor=west, minimum size=1in, inner xsep=5mm] at ($(current page.north west)+(1.3cm,-1.65cm)$) {\includegraphics[height=1.25cm]{Artifactory_logo_MARK_ORANGE.png}};
\end{tikzpicture}
}

\setlength{\headheight}{15pt}

% --------------------- Document Stuff
\pagestyle{fancy}
\begin{document}
\pagenumbering{gobble}

\begin{titlepage}

  % % DATE:
  % \date{\today}

  \centering % Centre everything on the title page

  %------------------------------------------------
  %	Top rules
  %------------------------------------------------

  \rule{\textwidth}{1pt} % Thick horizontal rule

  % \vspace{2pt}\vspace{-\baselineskip} % Whitespace between rules
  \vspace{2pt}\vspace{-10mm} % Whitespace between rules

  \rule{\textwidth}{0.4pt} % Thin horizontal rule

  \vspace{0.1\textheight} % Whitespace between the top rules and title

  %------------------------------------------------
  %	Title
  %------------------------------------------------

  {\Huge CONSTITUTION }\\[0.5\baselineskip] % Title line 1
  {\Large OF}\\[0.5\baselineskip] % Title line 2
  {\Huge PERTH ARTIFACTORY INC} % Title line 3

  \vspace{0.025\textheight} % Whitespace between the title and short horizontal rule

  \rule{0.4\textwidth}{0.4pt} % Short horizontal rule under the title

  \vspace{0.05\textheight}
  \today
  \vspace{0.2\textheight}

  % Logo 
  % \includegraphics[width=0.6\textwidth]{Artifactory_logo_COMBO-HORIZ_ORANGE.png}
  \includegraphics[width=0.4\textwidth]{Artifactory_logo_COMBO-ROUND_ORANGE.png}

  % \vspace{0.1\textheight}
  \vfill

  \rule{\textwidth}{0.4pt} % Thin horizontal rule
  \vspace{2pt}\vspace{-10mm} % Whitespace between rules
  \rule{\textwidth}{1pt} % Thick horizontal rule

\end{titlepage}

\pagenumbering{roman}

\tableofcontents

\pagebreak
\pagenumbering{arabic}
\subfile{sections/revision_history.tex}
% \pagebreak
% \subfile{sections/preface.tex}

\pagebreak
\subfile{sections/part_0.tex}
\pagebreak
\subfile{sections/part_1.tex}
\pagebreak
\subfile{sections/part_2.tex}
\pagebreak
\subfile{sections/part_3.tex}
\pagebreak
\subfile{sections/part_4.tex}
\pagebreak
\subfile{sections/part_5.tex}
\pagebreak
\subfile{sections/part_6.tex}
\pagebreak
\subfile{sections/part_7.tex}
\pagebreak
\subfile{sections/part_8.tex}
\pagebreak
\subfile{sections/own_rules_matters.tex}

\end{document}
